\rok{Фебруар 2026.}{А}

Други поправни тест из Алгоритама и структура података

Први практични тест одржан 19.02.2026.

\zadatak{1}{Insertion Sort}
Написати Insertion Sort алгоритам за сортирање целобројног низа дужине $n$.

\textbf{Додатне напомене за решавање задатка:}
\begin{itemize}
	\item Улаз је несортиран низ целих бројева.
	\item Излаз је сортирани низ целих бројева.
	\item У template-у се налази класа \texttt{RandomArrayGenerator} коју треба искористити за генерисање низа случајних бројева.
	\item У оквиру класе \texttt{InsertionSort} написати имплементацију за методу:
	\begin{Verbatim}[fontsize=\small]
public static void sort(long[] arr)
	\end{Verbatim}
	\item Обезбедити код у Main методи за тестирање алгоритма.
\end{itemize}

\zadatak{2}{Производ елемената између минималне вредности и последњег чвора}
Имплементирати методу која проналази прво појављивање минималне вредности у листи и враћа производ свих елемената од тог чвора до краја листе.

\textbf{Додатни захтеви за решавање задатка:}
\begin{itemize}
	\item Ако минимална вредност не постоји (празна листа), вратити 0.
	\item Алгоритам писати у оквиру класе \texttt{LinkedList}.
	\item Алгоритам мора имати временску комплексност $O(n)$.
	\item Захтеван потпис методе:
	\begin{Verbatim}[fontsize=\small]
public int productFromFirstMin()
	\end{Verbatim}
\end{itemize}

\textbf{Пример:}
\begin{itemize}
	\item \textbf{Улаз:} $4 \rightarrow 2 \rightarrow 3 \rightarrow 4$ \\
	      \textbf{Излаз:} $24$ ($2 * 3 * 4$)
	\item \textbf{Улаз:} $2 \rightarrow 3 \rightarrow 4 \rightarrow 4$ \\
	      \textbf{Излаз:} $96$ ($2 * 3 * 4 * 4$)
\end{itemize}

\zadatak{3}{Фузија бројева (The Number Fusion)}
Замисли да имаш низ бројева који улазе у експериментални реактор један по један. Сваки број представља енергетски ниво честице. Када нова честица падне на ону испод ње, покреће се процес фузије према следећим правилима:

\begin{itemize}
	\item \textbf{Фузија (спајање):} Ако је нова честица већа или једнака честици на врху, оне се спајају у једну нову честицу чија је вредност једнака њиховом збиру. Новонастала честица затим покушава да се споји са следећом честицом испод ње по истом правилу.
	\item \textbf{Одбијање:} Ако је нова честица (или новонастали збир) строго мања од честице на врху, фузија није могућа и она једноставно остаје да лежи на њој.
	\item \textbf{Празан реактор:} Ако је реактор празан, честица само заузима место на дну.
\end{itemize}

Твој циљ: Имплементирати методу која прима низ вредности честица и враћа финално стање у реактору (од дна ка врху).

\textbf{Додатни захтеви за решавање задатка:}
\begin{itemize}
	\item Временска сложеност: $O(n)$.
	\item Резултат мора бити низ у оригиналном редоследу (од честице на дну до оне на врху).
\end{itemize}

\textbf{Пример:}
\begin{itemize}
	\item \textbf{Улаз 1} (несортирани низ целих бројева): \texttt{[3, 2, 5]} \\
	      \textbf{Излаз 1:} \texttt{[10]} \\
	      ($5$ се спаја са $2$ у $7$; $7$ се спаја са $3$ у $10$)
	\item \textbf{Улаз 2} (несортирани низ целих бројева): \texttt{[10, 2, 1, 4]} \\
	      \textbf{Излаз 2:} \texttt{[10, 7]} \\
	      ($4$ се спаја са $1$ у $5$; $5$ се спаја са $2$ у $7$; $7$ је мање од $10$ и остаје)
\end{itemize}

\sablon{Шаблон другог поправног теста фебруар 2026. група А}{2025/Februar/GrupaA/AISP2025_PopravniTest2_Test1_GrupaA.linq}
